\section{Deterministic PDE and Capacity Proofs}
\label{app:deterministic}

This appendix proves \cref{thm:capacityscale,thm:localtd} without importing a
small-hole expansion.  Constants denoted by $C$ may change from line to line,
but are independent of $z\in Z$ and sufficiently small $\eps$.  We write
$K=K_{\eps,z}$, $\Omega_\eps=\Omega_{\eps,z}$,
$q_\eps=q_{\eps,z}$, and $\rho_\eps=\rho_z(\eps)$ when no confusion can
arise.  The background facts used below---trace and compact embedding,
Lax--Milgram, quasi-continuous representatives, positive distributions, and
local Green-kernel bounds---are standard results of Sobolev and elliptic
potential theory~\cite{evans2010pde,adams1996function,gilbarg2001elliptic}; all
problem-specific consequences are derived explicitly.

\subsection{Why regular shape calculus cannot nucleate the cavity}
\label{app:notshapepath}

\begin{proposition}[Cavity nucleation is not a regular shape path]
\label{prop:notshapepath}
Assume additionally that $\partial\Omega$ and $\partial B$ are connected.
Let $V\in C^1_c(\R^d;\R^d)$ and
$T_t=\operatorname{Id}+tV$.  For all sufficiently small $|t|$, $T_t$ is a
$C^1$ diffeomorphism, $T_t(\Omega)$ is diffeomorphic to $\Omega$, and
$\partial T_t(\Omega)=T_t(\partial\Omega)$ is connected.  Hence no such
regular perturbation satisfies
\begin{equation}
 T_t(\Omega)=\Omega\setminus\overline{(z+\eps B)}
 \label{eq:noshapepath}
\end{equation}
for any $\eps>0$ with $z+\eps\overline B\Subset\Omega$, because the right-hand
side has the two disjoint boundary components $\partial\Omega$ and
$z+\eps\partial B$.  Consequently \cref{eq:tddefinition} is not a directional
shape derivative along a regular deformation of the unperforated domain.
\end{proposition}

\begin{proof}[Proof of \cref{prop:notshapepath}]
We separate the analytic and topological parts.

\textbf{Step 1: $T_t$ is a local diffeomorphism for small $t$.}
Because $V\in C^1_c(\R^d;\R^d)$, the quantity
$L:=\|DV\|_{L^\infty(\R^d)}$ is finite.  Choose $|t|L<1$.  Then
\begin{equation}
 DT_t(x)=I+tDV(x)
\end{equation}
is invertible at every $x$: indeed,
$\|tDV(x)\|<1$, so the Neumann series
$\sum_{k=0}^{\infty}(-tDV(x))^k$ converges to $DT_t(x)^{-1}$.
The inverse-function theorem therefore makes $T_t$ a local $C^1$
diffeomorphism.

\textbf{Step 2: $T_t$ is globally injective and proper.}
For any $x,y\in\R^d$, the mean-value estimate gives
\begin{align}
 \|T_t(x)-T_t(y)\|
 &\ge \|x-y\|-|t|\,\|V(x)-V(y)\|\\
 &\ge (1-|t|L)\|x-y\|.
 \label{eq:bilipschitzlower}
\end{align}
Thus $T_t(x)=T_t(y)$ implies $x=y$.  Since $V$ has compact support,
$T_t(x)=x$ outside a compact set; hence $T_t$ is proper.  A proper injective
local diffeomorphism has open and closed image in connected $\R^d$, so its
image is all of $\R^d$.  Therefore $T_t$ is a global $C^1$ diffeomorphism.

\textbf{Step 3: a diffeomorphism preserves the boundary.}
A homeomorphism preserves closure and interior.  Hence
\begin{align}
 \partial T_t(\Omega)
 &=\overline{T_t(\Omega)}\setminus\operatorname{int}T_t(\Omega)\\
 &=T_t(\overline\Omega)\setminus T_t(\operatorname{int}\Omega)
 =T_t(\partial\Omega).
 \label{eq:boundarymapped}
\end{align}
Because $\partial\Omega$ is connected and $T_t$ is continuous,
$T_t(\partial\Omega)$ is connected.

\textbf{Step 4: an interior cavity creates a second boundary component.}
Put $K_{\eps,z}=z+\eps\overline B\Subset\Omega$ and
$\Omega_{\eps,z}=\Omega\setminus K_{\eps,z}$.  Positive separation of
$K_{\eps,z}$ and $\partial\Omega$ implies
\begin{equation}
 \partial\Omega_{\eps,z}
 =\partial\Omega\,\dot\cup\,(z+\eps\partial B).
 \label{eq:perforatedboundary}
\end{equation}
Both sets are nonempty, closed in the relative boundary, and separated by a
positive distance.  Since each is connected by assumption, they are exactly
two distinct connected components.  Therefore
$\partial\Omega_{\eps,z}$ is disconnected.

\textbf{Step 5: contradiction.}
If \cref{eq:noshapepath} held, \cref{eq:boundarymapped} would make the boundary
of its left-hand side connected, whereas \cref{eq:perforatedboundary} makes
the boundary of its right-hand side disconnected.  This is impossible.
Regular Eulerian shape calculus differentiates $J(T_t(\Omega))$ at $t=0$;
the perforated family does not lie on any such path.  Its capacity-normalized
limit is therefore a genuinely singular topological derivative rather than a
directional derivative at the original unperforated domain.  The argument
does not rule out additional asymptotic analysis of shape derivatives on
domains with pre-existing positive-radius cavities.
\end{proof}

\subsection{Functional-analytic foundation}

Equip $V=H^1_{\Gamma_D}(\Omega)$ with
\begin{equation}
 \|v\|_a^2:=a_\Omega(v,v).
 \label{eq:energynorm}
\end{equation}

\begin{proposition}[Boundedness and coercivity]
\label{prop:coercivity}
There are constants $c_a,C_a>0$ such that for all $v,w\in V$,
\begin{equation}
 |a_\Omega(v,w)|\le C_a\|v\|_{H^1}\|w\|_{H^1},
 \qquad
 a_\Omega(v,v)\ge c_a\|v\|_{H^1}^2.
 \label{eq:coercivity}
\end{equation}
The same coercivity constant applies to every perforated test space after
zero extension across the cavity.
\end{proposition}

\begin{proof}
\textbf{Step 1: boundedness.}
H\"older's inequality and the coefficient bounds give
\begin{align}
 |a_\Omega(v,w)|
 &\le \kappa_+\|\nabla v\|_2\|\nabla w\|_2
   +\|\sigma\|_\infty\|v\|_2\|w\|_2\\
 &\le C_a\|v\|_{H^1}\|w\|_{H^1}.
\end{align}

\textbf{Step 2: a Poincar\'e inequality.}
Because the trace vanishes on a boundary subset of positive
$(d-1)$-dimensional measure, there is $C_P<\infty$ such that
\begin{equation}
 \|v\|_{L^2(\Omega)}\le C_P\|\nabla v\|_{L^2(\Omega)}
 \qquad(v\in V).
 \label{eq:mixedpoincare}
\end{equation}
Indeed, if this failed, a normalized sequence with unit $L^2$ norm and
vanishing gradients would converge in $L^2$ to a nonzero constant; continuity
of the trace would force that constant to vanish on $\Gamma_D$, a
contradiction.

\textbf{Step 3: coercivity.}
Since $\sigma\ge0$ and $\kappa\ge\kappa_-$,
\begin{equation}
 a_\Omega(v,v)\ge\kappa_-\|\nabla v\|_2^2
 \ge\frac{\kappa_-}{1+C_P^2}\|v\|_{H^1}^2.
\end{equation}
If $v$ is a perforated-domain test function, its zero extension across $K$
belongs to $V$ because its trace on $\partial K$ is zero.  Applying the same
inequality to that extension proves the final assertion.
\end{proof}

\begin{proposition}[Well-posed state, adjoint, and perforated state]
\label{prop:wellposedness}
The weak problems in \cref{eq:stateweak,eq:adjointweak} have unique solutions.
For every sufficiently small $\eps$, the perforated state with trace $g$ on
$\Gamma_D$ and $h_z$ on $\partial K$ also exists uniquely.  Its extension by
$h_z$ inside $K$ belongs to $H^1(\Omega)$.
\end{proposition}

\begin{proof}
\textbf{Step 1: the state.}
Choose the fixed lift $g\in H^1(\Omega)$ and write $u=g+u_0$ with $u_0\in V$.
Then
\begin{equation}
 a_\Omega(u_0,v)=\ell_\Omega(v)-a_\Omega(g,v)
 \qquad(v\in V).
\end{equation}
The right side is a bounded linear functional.  Boundedness and coercivity
from \cref{prop:coercivity} allow the Lax--Milgram theorem to give a unique
$u_0$, hence a unique $u$.

\textbf{Step 2: the adjoint.}
Because $j'(u)$ is bounded on $L^2(\Omega)$, its restriction to $V$ is bounded
by the continuous embedding $V\hookrightarrow L^2$.  Another application of
Lax--Milgram gives a unique $p\in V$ satisfying
\cref{eq:adjointweak}.

\textbf{Step 3: the perforated state.}
Define
\begin{equation}
 V_\eps:=\{v\in H^1(\Omega_\eps):
 \operatorname{Tr}v=0\text{ on }\Gamma_D\cup\partial K\}.
\end{equation}
Since $K\Subset\Omega$, a smooth cutoff supplies an $H^1(\Omega_\eps)$ lift
with outer trace $g$ and inner trace $h_z$.  Subtracting that lift reduces the
problem to $V_\eps$.  The zero-extension argument in
\cref{prop:coercivity} supplies uniform coercivity, so Lax--Milgram again gives
existence and uniqueness.

\textbf{Step 4: extension across the cavity.}
The trace of $u_{\eps,z}$ on $\partial K$ equals the trace of the constant
function $h_z$ inside $K$.  The Sobolev gluing lemma therefore shows that the
piecewise definition $\widetilde u_{\eps,z}=u_{\eps,z}$ on $\Omega_\eps$ and
$\widetilde u_{\eps,z}=h_z$ on $K$ belongs to $H^1(\Omega)$.
\end{proof}

\subsection{Equilibrium potential and capacitary measure}

Let
\begin{equation}
 \mathcal A_{\eps,z}:=\{v\in V:\widetilde v\ge1
 \text{ quasi-everywhere on }K_{\eps,z}\}.
\end{equation}

\begin{lemma}[Equilibrium potential]
\label{lem:equilibrium}
The class $\mathcal A_{\eps,z}$ is nonempty, closed, and convex.  It has a
unique minimum-energy element $q_\eps$.  This minimizer satisfies
\begin{align}
 0\le q_\eps&\le1\quad\text{quasi-everywhere},\notag\\
 q_\eps&=1\quad\text{quasi-everywhere on }K,
 \label{eq:qrange}\\
 a_\Omega(q_\eps,\varphi)&=0
 \quad\text{if }\varphi\in V
 \text{ and }\widetilde\varphi=0\text{ q.e. on }K.
 \label{eq:qharmonic}
\end{align}
\end{lemma}

\begin{proof}
\textbf{Step 1: nonemptiness and convexity.}
Because $K\Subset\Omega$, a compactly supported smooth cutoff equal to one on
a neighborhood of $K$ belongs to $\mathcal A_{\eps,z}$.  Convexity is
immediate from the pointwise inequality.

\textbf{Step 2: weak closedness.}
Let $v_n\in\mathcal A_{\eps,z}$ converge weakly in $V$ to $v$.  Mazur's lemma
provides convex combinations of tails of $(v_n)$ that converge strongly to
$v$ in $V$.  A subsequence of a strongly $H^1$-convergent sequence converges
quasi-everywhere.  Every convex combination remains at least one
quasi-everywhere on $K$, hence so does $v$.  Thus $v\in\mathcal A_{\eps,z}$.

\textbf{Step 3: existence and uniqueness.}
A minimizing sequence is bounded by coercivity.  Weak compactness and the
preceding closure give a weak limit in $\mathcal A_{\eps,z}$; weak lower
semicontinuity of the squared Hilbert norm gives a minimizer.  Strict
convexity of $a_\Omega(v,v)$ on $V$ gives uniqueness.

\textbf{Step 4: truncation.}
The map $T(s)=\min(1,\max(0,s))$ is one-Lipschitz, satisfies
$|T(s)|\le|s|$, and preserves the constraint on $K$.  The Sobolev chain rule
therefore gives
\begin{equation}
 a_\Omega(T(q_\eps),T(q_\eps))
 \le a_\Omega(q_\eps,q_\eps).
\end{equation}
Uniqueness forces $T(q_\eps)=q_\eps$, proving \cref{eq:qrange}.

\textbf{Step 5: harmonicity away from the obstacle.}
If $\widetilde\varphi=0$ quasi-everywhere on $K$, then
$q_\eps+t\varphi$ is admissible for every $t\in\R$.  Differentiating its
quadratic energy at $t=0$ in both signs gives
$a_\Omega(q_\eps,\varphi)=0$.
\end{proof}

\begin{lemma}[Capacitary measure and localization]
\label{lem:capacitymeasure}
There is a finite positive Radon measure $\mu_\eps$ supported on $\overline K$
such that
\begin{equation}
 a_\Omega(q_\eps,v)=\int_{\overline K}\widetilde v\,\mathrm d\mu_\eps
 \quad(v\in V),
 \qquad
 \mu_\eps(\overline K)=\rho_\eps.
 \label{eq:measureidentity}
\end{equation}
Consequently, for every $v\in V\cap C^1(U)$,
\begin{equation}
 |a_\Omega(q_\eps,v)-\rho_\eps v(z)|
 \le \eps R_B\rho_\eps\|\nabla v\|_{L^\infty(U)},
 \label{eq:capacitylocalization}
\end{equation}
where $B\subset B_{R_B}(0)$.
\end{lemma}

\begin{proof}
\textbf{Step 1: positivity.}
For every nonnegative $\varphi\in C_c^\infty(\Omega)$ and $t>0$,
$q_\eps+t\varphi$ is admissible.  Minimality gives
\begin{equation}
 0\le\frac{a(q_\eps+t\varphi,q_\eps+t\varphi)
               -a(q_\eps,q_\eps)}{2t}
 =a(q_\eps,\varphi)+\frac t2a(\varphi,\varphi).
\end{equation}
Letting $t\downarrow0$ yields $a(q_\eps,\varphi)\ge0$.  Thus
$\varphi\mapsto a(q_\eps,\varphi)$ is a positive distribution and hence a
positive Radon measure $\mu_\eps$.

\textbf{Step 2: support.}
By \cref{eq:qharmonic}, this distribution vanishes on every test function
supported in $\Omega\setminus\overline K$.  Therefore
$\operatorname{supp}\mu_\eps\subseteq\overline K$.

\textbf{Step 3: obtain the mass without extending the measure identity.}
Choose $\chi\in C_c^\infty(U)$ with $\chi=1$ on a neighborhood of $K$.
The function $q_\eps-\chi$ belongs to $V$ and vanishes quasi-everywhere
on $K$. Thus \cref{eq:qharmonic}, applied directly in $V$, gives
\[
 \rho_\eps=a(q_\eps,q_\eps)=a(q_\eps,\chi)
 =\int\chi\,d\mu_\eps=\mu_\eps(K).
\]
The integral identity here is used only for the smooth compactly supported
function $\chi$, for which Steps 1--2 already prove it.

\textbf{Step 4: the measure does not charge polar sets.}
For a compact capacity-zero set $F\subset K$, the smooth-cutoff
characterization of Sobolev capacity gives nonnegative
$\eta_n\in C_c^\infty(\Omega)$ that are at least one near $F$ and satisfy
$\|\eta_n\|_a\to0$. This interior characterization agrees with the relative
capacity: multiply admissible approximations by a fixed interior cutoff,
use norm equivalence, and the usual smooth approximation characterization
of Sobolev capacity. Only compact subsets of $K\Subset\Omega$ are needed.
Positivity and the already established smooth-test identity imply
\[
 \mu_\eps(F)\le\int\eta_n\,d\mu_\eps
 =a(q_\eps,\eta_n)
 \le\rho_\eps^{1/2}\|\eta_n\|_a\longrightarrow0.
\]
Inner regularity of the finite Radon measure then proves the same assertion
for every Borel capacity-zero subset of $K$.

\textbf{Step 5: extend locally by an energy bound.}
For $\varphi\in C_c^\infty(\Omega)$, use the smooth nonnegative functions
$T_\delta(\varphi)=\sqrt{\varphi^2+\delta^2}-\delta$.
They have the same compact support, converge to $|\varphi|$, and satisfy
$\|T_\delta(\varphi)\|_a\le\|\varphi\|_a$ by the chain rule and $\sigma\ge0$.
Consequently, dominated convergence on the finite measure space gives
\[
 \int|\varphi|\,d\mu_\eps
 =\lim_{\delta\downarrow0}a(q_\eps,T_\delta(\varphi))
 \le\rho_\eps^{1/2}\|\varphi\|_a.
\]
If smooth compactly supported $\varphi_n$ converge to $w\in H^1_0(\Omega)$,
this inequality makes them Cauchy in $L^1(\mu_\eps)$. A subsequence converges
quasi-everywhere to $\widetilde w$ and, by Step 4, also $\mu_\eps$-almost
everywhere. Its $L^1(\mu_\eps)$ limit is therefore $\widetilde w$.
Passing to the limit in the smooth-test identity proves
$a(q_\eps,w)=\int\widetilde w\,d\mu_\eps$ for $w\in H^1_0(\Omega)$.

\textbf{Step 6: extend to the mixed-boundary space by localization.}
For an arbitrary $v\in V$, $\chi v\in H^1_0(\Omega)$, while
$v-\chi v$ vanishes on $K$. Orthogonality in \cref{eq:qharmonic} and Step 5 give
\[
 a(q_\eps,v)=a(q_\eps,\chi v)
 =\int_K\chi\widetilde v\,d\mu_\eps
 =\int_K\widetilde v\,d\mu_\eps.
\]
This avoids assuming that $C_c^\infty(\Omega)$ is dense in $V$, which would
be false when nonzero traces are allowed on $\Gamma_N$.

\textbf{Step 7: localization of a smooth observable.}
For $v\in V\cap C^1(U)$, every $x\in K$ obeys $|x-z|\le\eps R_B$.
Using the mass from Step 3,
\begin{align}
 |a(q_\eps,v)-\rho_\eps v(z)|
 &=\left|\int_K(v(x)-v(z))\,d\mu_\eps(x)\right|\\
 &\le\eps R_B\|\nabla v\|_{L^\infty(U)}\rho_\eps.
\end{align}
This proves \cref{eq:capacitylocalization} and completes the extension argument.
\end{proof}

\subsection{Capacity scale}

\begin{proof}[Proof of \cref{thm:capacityscale}]
Existence and uniqueness follow from \cref{lem:equilibrium}; coercivity and
the fact that $K$ contains a ball imply $\rho_\eps>0$.

\paragraph{The three-dimensional upper bound.}
Let $P_B$ be the Newtonian equilibrium potential of $B$:
$P_B=1$ quasi-everywhere on $B$, $P_B\to0$ at infinity, and
\begin{equation}
 \int_{\R^3}|\nabla P_B|^2\,\mathrm dy=\operatorname{cap}(B).
 \label{eq:newtoniancapacity}
\end{equation}
Choose a fixed cutoff $\chi\in C_c^\infty(\Omega)$ that equals one on a
neighborhood of $Z$.  The function
\begin{equation}
 v_{\eps,z}(x):=\chi(x)P_B((x-z)/\eps)
\end{equation}
is admissible for all sufficiently small $\eps$.  After the change of
variables $x=z+\eps y$, local uniform convergence
$\kappa(z+\eps y)\to\kappa(z)$ and the finite energy of $P_B$ yield
\begin{equation}
 \int_\Omega\kappa|\nabla v_{\eps,z}|^2\,\mathrm dx
 =\eps\kappa(z)\operatorname{cap}(B)+o(\eps).
 \label{eq:capupper3dgrad}
\end{equation}
The terms in which the gradient hits $\chi$ occur a fixed positive distance
from $z$, where $P_B((x-z)/\eps)=O(\eps)$; they are $O(\eps^2)$.  Likewise
$\int\sigma v_{\eps,z}^2=O(\eps^2)=o(\eps)$.  Minimality of $q_\eps$ gives
\begin{equation}
 \limsup_{\eps\downarrow0}\frac{\rho_\eps}{\eps}
 \le\kappa(z)\operatorname{cap}(B).
 \label{eq:capupper3d}
\end{equation}

\paragraph{The three-dimensional lower bound.}
Define $Q_\eps(y)=q_\eps(z+\eps y)$ on
$(\Omega-z)/\eps$. The upper bound makes $\int|\nabla Q_\eps|^2$ uniformly bounded. Moreover, Sobolev embedding and coercivity on the fixed domain give $\|q_\eps\|_{L^6(\Omega)}\le C\rho_\eps^{1/2}=O(\eps^{1/2})$. Rescaling yields a uniform $L^6$ bound for $Q_\eps$. Local weak compactness and a diagonal subsequence therefore give a limit $Q$ with $Q\in L^6(\R^3)$ and $\nabla Q\in L^2(\R^3)$, hence $Q\in D^{1,2}(\R^3)$. The cavity constraint passes to the limit by weak closedness on any fixed ball containing $B$, so $Q\ge1$ quasi-everywhere on $B$.  For every fixed $R$,
uniform convergence of the coefficient and weak lower semicontinuity give
\begin{align}
 \liminf_{\eps\downarrow0}\frac{\rho_\eps}{\eps}
 &\ge\liminf_{\eps\downarrow0}
 \int_{B_R}\kappa(z+\eps y)|\nabla Q_\eps|^2\,\mathrm dy\\
 &\ge\kappa(z)\int_{B_R}|\nabla Q|^2\,\mathrm dy.
\end{align}
Letting $R\to\infty$ and then using the variational definition in
\cref{eq:newtoniancapacity} yields
\begin{equation}
 \liminf_{\eps\downarrow0}\frac{\rho_\eps}{\eps}
 \ge\kappa(z)\operatorname{cap}(B).
 \label{eq:caplower3d}
\end{equation}
Together, \cref{eq:capupper3d,eq:caplower3d} prove the $d=3$ formula.

\paragraph{The two-dimensional upper bound.}
By compactness, choose $r_0>0$ so that
\begin{equation}
 B_{2r_0}(z)\Subset\Omega\qquad\text{for every }z\in Z.
\end{equation}
Let $B\subset B_{R_B}(0)$.  Put
$L_\eps=\log(r_0/(\eps R_B))$ and use the radial competitor
\begin{equation}
 v(r)=
 \begin{cases}
 1,&r\le\eps R_B,\\
 \log(r_0/r)/L_\eps,&\eps R_B<r<r_0,\\
 0,&r\ge r_0.
 \end{cases}
\end{equation}
A direct polar-coordinate calculation gives
\begin{equation}
 \int_{\eps R_B}^{r_0}|v'(r)|^2\,2\pi r\,\mathrm dr
 =\frac{2\pi}{L_\eps}.
 \label{eq:annulusenergy}
\end{equation}
Uniform continuity of $\kappa$ and then $r_0\downarrow0$ show
\begin{equation}
 \limsup_{\eps\downarrow0}
 |\log\eps|\rho_\eps\le2\pi\kappa(z).
 \label{eq:capupper2d}
\end{equation}
The zero-order energy is $O(L_\eps^{-2})+O(\eps^2)$ and therefore does not
alter this limit.

\paragraph{The two-dimensional lower bound.}
Choose $r_B>0$ with $B_{r_B}(0)\subset B$.  Fix $r_0$ as above and let
$\delta>0$ be such that
$\kappa(x)\ge\kappa(z)-\delta$ on $B_{r_0}(z)$.  For an admissible $v$, let
$c$ be its circular mean on $\partial B_{r_0}(z)$.  Applying Cauchy--Schwarz
to the radial derivative of the circular mean gives
\begin{equation}
 \int_{B_{r_0}(z)\setminus B_{\eps r_B}(z)}|\nabla v|^2\,\mathrm dx
 \ge\frac{2\pi(1-c)_+^2}{\log(r_0/(\eps r_B))}.
 \label{eq:annuluslower}
\end{equation}
On the fixed exterior domain $\Omega\setminus B_{r_0}(z)$, the trace theorem,
the Poincar\'e inequality on functions vanishing at $\Gamma_D$, and
\cref{eq:coercivity} give
\begin{equation}
 a_{\Omega\setminus B_{r_0}(z)}(v,v)\ge C_0c^2
 \label{eq:outerresistance}
\end{equation}
with $C_0>0$ uniform on $Z$.  Minimizing the sum of
$A(1-c)_+^2+C_0c^2$, where
$A=2\pi(\kappa(z)-\delta)/\log(r_0/(\eps r_B))$, gives
\begin{equation}
 \rho_\eps\ge\frac{AC_0}{A+C_0}.
\end{equation}
First let $\eps\downarrow0$ and then $\delta\downarrow0$ to obtain
\begin{equation}
 \liminf_{\eps\downarrow0}|\log\eps|\rho_\eps
 \ge2\pi\kappa(z).
 \label{eq:caplower2d}
\end{equation}
This and \cref{eq:capupper2d} prove the $d=2$ formula.

\paragraph{Uniformity and comparison with volume.}
All cutoffs, radii, trace constants, and coefficient moduli can be chosen
uniformly because $Z$ is compactly contained in $\Omega$.  A contradiction
argument applied to any sequence $z_n\in Z$, $\eps_n\downarrow0$ upgrades the
pointwise limits to uniform ones.  Finally,
$\eps^3=o(\eps)$ in $d=3$ and
$\eps^2=o(1/|\log\eps|)$ in $d=2$, proving
$\eps^d=o(\rho_\eps)$.
\end{proof}

\subsection{Potential and corrector estimates}

\begin{lemma}[$L^2$ size of the equilibrium potential]
\label{lem:qLtwo}
For $d\in\{2,3\}$,
\begin{equation}
 \|q_\eps\|_{L^2(\Omega)}\le C\rho_\eps.
 \label{eq:qLtwo}
\end{equation}
\end{lemma}

\begin{proof}
The coercive mixed-boundary operator has a positive Green kernel $G(x,y)$.
For $y$ in the fixed interior neighborhood $U$, its local singularity and
regular part satisfy
\begin{equation}
 0\le G(x,y)\le C
 \begin{cases}
 1+|\log|x-y||,&d=2,\\
 |x-y|^{-1}+1,&d=3.
 \end{cases}
 \label{eq:greenbound}
\end{equation}
These bounds follow by freezing the uniformly elliptic principal coefficient,
using the Laplace fundamental solution as a local parametrix, and applying the
maximum principle plus interior elliptic estimates to the regular remainder.
Both right sides are square-integrable in $x$ on a bounded domain, uniformly
for $y\in U$; hence
$\sup_{y\in U}\|G(\cdot,y)\|_2\le C_G$.

By \cref{eq:measureidentity}, $q_\eps$ is the unique weak solution of
$Lq_\eps=\mu_\eps$ with homogeneous outer conditions, so
\begin{equation}
 q_\eps(x)=\int_KG(x,y)\,\mathrm d\mu_\eps(y).
\end{equation}
Minkowski's integral inequality and
$\mu_\eps(K)=\rho_\eps$ yield
\begin{equation}
 \|q_\eps\|_2
 \le\int_K\|G(\cdot,y)\|_2\,\mathrm d\mu_\eps(y)
 \le C_G\rho_\eps.
\end{equation}
\end{proof}

\begin{lemma}[State corrector]
\label{lem:statecorrector}
Let $w_\eps:=\widetilde u_{\eps,z}-u$ and define
\begin{equation}
 r_\eps:=w_\eps+(u(z)-h_z)q_\eps.
 \label{eq:correctordefinition}
\end{equation}
Then
\begin{align}
 \|r_\eps\|_a&\le C\eps^{d/2},
 \label{eq:correctorenergy}\\
 |a_\Omega(r_\eps,p)|&\le C\eps^{d/2},
 \label{eq:correctoradjoint}\\
 \|w_\eps\|_{L^2(\Omega)}&\le C\rho_\eps.
 \label{eq:stateLtwo}
\end{align}
In particular, $a_\Omega(r_\eps,p)=o(\rho_\eps)$ uniformly on $Z$.
\end{lemma}

\begin{proof}
\textbf{Step 1: the corrector equation and trace.}
The state equations on $\Omega_\eps$ imply that $w_\eps$ is $a$-harmonic
against every test function vanishing on both the outer Dirichlet boundary and
$\partial K$.  The same is true of $q_\eps$ by
\cref{eq:qharmonic}.  Therefore $r_\eps$ is $a$-harmonic on
$\Omega_\eps$.  On the cavity boundary,
\begin{equation}
 r_\eps=h_z-u+(u(z)-h_z)=u(z)-u.
 \label{eq:correctortrace}
\end{equation}

\textbf{Step 2: a localized lift.}
Choose $R>R_B$ and $\chi\in C_c^\infty(B_{2R}(0))$ with
$\chi=1$ on $B_R(0)$.  For sufficiently small $\eps$, define
\begin{equation}
 \eta_\eps(x):=\chi((x-z)/\eps)(u(z)-u(x)).
\end{equation}
It has the trace in \cref{eq:correctortrace}.  Local $C^1$ regularity gives
$|u(z)-u(x)|\le C\eps$ on its support and
\begin{align}
 |\nabla\eta_\eps|
 &\le\eps^{-1}|\nabla\chi((x-z)/\eps)|
              \cdot|u(z)-u(x)|\notag\\
 &\quad+|\chi((x-z)/\eps)|\cdot|\nabla u(x)|\notag\\
 &\le C.
\end{align}
Because the support has volume $O(\eps^d)$,
\begin{equation}
 a_\Omega(\eta_\eps,\eta_\eps)\le C\eps^d.
 \label{eq:liftenergy}
\end{equation}

\textbf{Step 3: energy minimality.}
On $\Omega_\eps$, both $r_\eps$ and $\eta_\eps$ have the same boundary trace.
Thus $r_\eps-\eta_\eps$ is an admissible homogeneous test function, and
$a_{\Omega_\eps}(r_\eps,r_\eps-\eta_\eps)=0$.  Consequently
\begin{equation}
 a_{\Omega_\eps}(\eta_\eps,\eta_\eps)
 =a_{\Omega_\eps}(r_\eps,r_\eps)
  +a_{\Omega_\eps}(\eta_\eps-r_\eps,\eta_\eps-r_\eps).
\end{equation}
Inside $K$, the extension of $r_\eps$ equals $u(z)-u(x)$ and has energy
$O(\eps^d)$.  Combining this fact with \cref{eq:liftenergy} proves
\cref{eq:correctorenergy}.

\textbf{Step 4: adjoint pairing.}
Cauchy--Schwarz in the energy inner product gives
\begin{equation}
 |a(r_\eps,p)|\le\|r_\eps\|_a\|p\|_a
 \le C\eps^{d/2},
\end{equation}
which is \cref{eq:correctoradjoint}.

\textbf{Step 5: $L^2$ size of the state perturbation.}
From \cref{eq:correctordefinition}, Poincar\'e's inequality, and
\cref{lem:qLtwo},
\begin{align}
 \|w_\eps\|_2
 &\le |u(z)-h_z|\,\|q_\eps\|_2+\|r_\eps\|_2\\
 &\le C\rho_\eps+C\eps^{d/2}.
\end{align}
The capacity scales show $\eps^{d/2}=o(\rho_\eps)$ for $d=2,3$, proving
\cref{eq:stateLtwo}.  The same comparison applied to
\cref{eq:correctoradjoint} proves the final statement.
\end{proof}

\subsection{Objective expansion}

\begin{proof}[Proof of \cref{thm:localtd}]
Fix $z\in Z$ and abbreviate $a_z=u(z)-h_z$.

\textbf{Step 1: Taylor expansion of the state objective.}
By \cref{eq:stateLtwo}, $w_\eps\to0$ in $L^2$.  The integral form of the
second-order Fr\'echet remainder is
\begin{equation}
 j(u+w_\eps)-j(u)-j'(u)[w_\eps]
 =\int_0^1(1-t)j''(u+tw_\eps)[w_\eps,w_\eps]\,\mathrm dt.
\end{equation}
The uniform bound on $j''$ and \cref{eq:stateLtwo} imply
\begin{equation}
 |j(u+w_\eps)-j(u)-j'(u)[w_\eps]|
 \le C\rho_\eps^2.
 \label{eq:objectiveTaylorbound}
\end{equation}

\textbf{Step 2: explicit domain cost.}
The only explicit domain change is
\begin{equation}
 \left|\int_{\Omega_\eps}\psi-\int_\Omega\psi\right|
 =\left|\int_K\psi\right|
 \le\|\psi\|_\infty|B|\eps^d.
 \label{eq:domaincostbound}
\end{equation}

\textbf{Step 3: adjoint substitution.}
The extension $w_\eps$ has zero trace on $\Gamma_D$, hence belongs to $V$.
By \cref{eq:adjointweak,eq:correctordefinition},
\begin{equation}
 j'(u)[w_\eps]
 =a_\Omega(w_\eps,p)
 =-a_z a_\Omega(q_\eps,p)+a_\Omega(r_\eps,p).
 \label{eq:adjointcorrector}
\end{equation}

\textbf{Step 4: localization of the capacitary charge.}
Applying \cref{eq:capacitylocalization} to the locally $C^1$ adjoint gives
\begin{equation}
 a_\Omega(q_\eps,p)
 =\rho_\eps p(z)+e_{\eps,z},
 \qquad |e_{\eps,z}|\le C\eps\rho_\eps.
 \label{eq:adjointlocalization}
\end{equation}
Combining \cref{eq:correctoradjoint,eq:adjointcorrector,eq:adjointlocalization}
yields
\begin{equation}
 j'(u)[w_\eps]
 =-\rho_\eps(u(z)-h_z)p(z)
  +O(\eps\rho_\eps)+O(\eps^{d/2}).
 \label{eq:firstvariationfinal}
\end{equation}

\textbf{Step 5: collect and express the remainder in capacity units.}
Equations \cref{eq:objectiveTaylorbound,eq:domaincostbound,eq:firstvariationfinal}
give
\begin{align}
 &\left|\J(\Omega_{\eps,z})-\J(\Omega)
 +\rho_\eps(u(z)-h_z)p(z)\right|\notag\\
 &\hspace{16mm}\le
 C\left(\rho_\eps^2+\eps\rho_\eps
        +\eps^{d/2}+\eps^d\right).
 \label{eq:rawremainder}
\end{align}
For $d=3$, $\rho_\eps\asymp\eps$, so every term on the right is bounded by
$C\rho_\eps^{3/2}$.  For $d=2$,
$\rho_\eps\asymp1/|\log\eps|$, and the elementary limit
$\eps|\log\eps|^m\to0$ for every fixed $m$ shows that every term is bounded by
$C\rho_\eps^2$.  This proves \cref{eq:quantitativesmallhole} with
$\alpha_3=1/2$ and $\alpha_2=1$.

\textbf{Step 6: identify the derivative.}
Divide \cref{eq:rawremainder} by $\rho_\eps>0$ and let $\eps\downarrow0$.
The normalized remainder vanishes, so the limit in
\cref{eq:tddefinition} exists and equals
$-(u(z)-h_z)p(z)$.  Every constant above is uniform on $Z$, proving uniform
convergence and \cref{eq:localproduct}.
\end{proof}
\section{Cross-Fitted Estimator Proofs}
\label{app:ustat}

Fix a candidate $z$.  Write
\begin{equation}
 A_i:=U_i-h_z,\qquad B_i:=P_i,
 \qquad \mu_A:=\E A_i,\quad\mu_B:=\E B_i,
\end{equation}
and $X_i=(A_i,B_i)$.  Under the exact-oracle assumption,
$\theta:=-\mu_A\mu_B=\DT(z)$.

\subsection{Algebra and unbiasedness}

\begin{lemma}[Linear-time identity]
\label{lem:lineartime}
For every $N\ge2$,
\begin{equation}
 \sum_{i\ne j}A_iB_j
 =\left(\sum_{i=1}^NA_i\right)
  \left(\sum_{j=1}^NB_j\right)
  -\sum_{i=1}^NA_iB_i.
 \label{eq:lineartimeproof}
\end{equation}
\end{lemma}

\begin{proof}
Distributivity gives
$(\sum_iA_i)(\sum_jB_j)=\sum_{i,j}A_iB_j$.  The index set of this double sum
is the disjoint union of the diagonal $\{i=j\}$ and the off-diagonal
$\{i\ne j\}$.  Subtracting the diagonal proves the identity.  Thus the
estimator needs three scalar sums and $O(N)$ arithmetic.
\end{proof}

\begin{proof}[Proof of \cref{thm:unbiased}]
\textbf{Step 1: factor only across different sample pairs.}
For $i\ne j$, $X_i$ and $X_j$ are independent.  Therefore $A_i$ and $B_j$ are
independent even though $A_i$ and $B_i$ need not be, and
\begin{equation}
 \E[-A_iB_j]=-(\E A_i)(\E B_j)=-\mu_A\mu_B=\theta.
 \label{eq:crossmean}
\end{equation}

\textbf{Step 2: average the ordered terms.}
There are exactly $N(N-1)$ ordered pairs $(i,j)$ with $i\ne j$.  Linearity of
expectation and \cref{eq:crossmean} give
\begin{equation}
 \E\widehat{\DT}_N
 =\frac{1}{N(N-1)}\sum_{i\ne j}\theta=\theta.
\end{equation}
Finally, \cref{eq:oraclemeans,eq:localproduct} imply
$\theta=-(u(z)-h_z)p(z)=\DT(z)$.  At no point was
$\E[A_iB_i]$ factored, so within-pair dependence is unrestricted.
\end{proof}

\subsection{Exact variance decomposition}

Let $A_i^\circ=A_i-\mu_A$ and $B_i^\circ=B_i-\mu_B$.  Recall the symmetric
kernel
\begin{equation}
 H(X_i,X_j)=-\frac12(A_iB_j+A_jB_i),
\end{equation}
for which
\begin{equation}
 \widehat{\DT}_N
 =\binom N2^{-1}\sum_{1\le i<j\le N}H(X_i,X_j).
 \label{eq:ustatunordered}
\end{equation}

\begin{lemma}[Hoeffding projections]
\label{lem:projections}
The first-order and degenerate projections are
\begin{align}
 h_1(X_i)
 &:=-\frac12\left(\mu_BA_i^\circ+\mu_AB_i^\circ\right),
 \label{eq:hprojectionone}\\
 h_2(X_i,X_j)
 &:=-\frac12\left(A_i^\circ B_j^\circ
                  +A_j^\circ B_i^\circ\right).
 \label{eq:hprojectiontwo}
\end{align}
They obey
\begin{align}
 \E h_1(X_i)&=0,\notag\\
 \E[h_2(X_i,X_j)\mid X_i]&=0,\notag\\
 \E[h_2(X_i,X_j)\mid X_j]&=0.
 \label{eq:degeneracy}
\end{align}
\end{lemma}

\begin{proof}
Because $X_j$ is independent of $X_i$,
\begin{align}
 \E[H(X_i,X_j)\mid X_i]
 &=-\frac12(A_i\mu_B+B_i\mu_A).
\end{align}
Subtracting $\theta=-\mu_A\mu_B$ and collecting centered terms gives
\cref{eq:hprojectionone}.  Next expand
\begin{align}
 H(X_i,X_j)
 &=-\mu_A\mu_B
 -\frac12\{\mu_B(A_i^\circ+A_j^\circ)
            +\mu_A(B_i^\circ+B_j^\circ)\}\notag\\
 &\quad-\frac12(A_i^\circ B_j^\circ+A_j^\circ B_i^\circ).
\end{align}
Subtracting $\theta+h_1(X_i)+h_1(X_j)$ proves
\cref{eq:hprojectiontwo}.  Centering and independence across $i\ne j$ give
all three identities in \cref{eq:degeneracy}.
\end{proof}

\begin{proof}[Proof of \cref{thm:variance}]
\textbf{Step 1: decompose the U-statistic.}
Substitute
$H=\theta+h_1(X_i)+h_1(X_j)+h_2(X_i,X_j)$ into
\cref{eq:ustatunordered}.  Every $h_1(X_i)$ appears in $N-1$ unordered pairs,
and $\binom N2=N(N-1)/2$, hence
\begin{equation}
 \widehat{\DT}_N-\theta
 =\frac2N\sum_{i=1}^Nh_1(X_i)
 +\binom N2^{-1}\sum_{i<j}h_2(X_i,X_j).
 \label{eq:hoeffdingdecomposition}
\end{equation}

\textbf{Step 2: prove orthogonality.}
For any $k$ and pair $i<j$, the covariance
$\E[h_1(X_k)h_2(X_i,X_j)]$ is zero.  If $k$ is distinct from both indices,
independence and centering suffice.  If $k=i$, condition on $X_i$ and use
$\E[h_2\mid X_i]=0$; the case $k=j$ is identical.

Two different degenerate terms are also uncorrelated.  Disjoint pairs are
independent.  If $(i,j)$ and $(i,k)$ share one index, then conditioning on
$X_i$ and using independence of $X_j,X_k$ gives zero by degeneracy.  Thus only
identical unordered pairs contribute to the variance of the second sum.

\textbf{Step 3: take variances.}
With $\zeta_1=\Var(h_1(X_1))$ and
$\zeta_2=\Var(h_2(X_1,X_2))$, orthogonality gives
\begin{align}
 \Var(\widehat{\DT}_N)
 &=\frac{4}{N^2}N\zeta_1
 +\binom N2^{-2}\binom N2\zeta_2\\
 &=\frac4N\zeta_1+\frac{2}{N(N-1)}\zeta_2.
 \label{eq:variancefromprojection}
\end{align}

\textbf{Step 4: compute $\zeta_1$.}
Squaring \cref{eq:hprojectionone} and taking expectations gives
\begin{equation}
 \zeta_1=\frac14\left(
 \mu_B^2\Var A+\mu_A^2\Var B
 +2\mu_A\mu_B\Cov(A,B)\right),
\end{equation}
which is \cref{eq:zetaone}.

\textbf{Step 5: compute $\zeta_2$.}
Let $c_{AB}=\E[A_1^\circ B_1^\circ]$.  From
\cref{eq:hprojectiontwo}, independence across sample pairs yields
\begin{align}
 \E h_2^2
 &=\frac14\E\left[
 (A_1^\circ B_2^\circ)^2+(A_2^\circ B_1^\circ)^2
 +2A_1^\circ B_1^\circ A_2^\circ B_2^\circ\right]\\
 &=\frac14\left(2\Var A\,\Var B+2c_{AB}^2\right),
\end{align}
which is \cref{eq:zetatwo}.  The second-moment assumptions make all quantities
finite.  Since $A,B$, and $H$ are evaluated on the unperforated domain, no
$\eps$ occurs anywhere in \cref{eq:variancefromprojection}.
\end{proof}

\subsection{Finite-sample concentration}

\begin{proof}[Proof of \cref{thm:concentration}]
Let $m=\lfloor N/2\rfloor$ and let $\Pi$ be a uniformly random permutation of
$\{1,\ldots,N\}$, independent of the data.  Define the disjoint-pair average
\begin{equation}
 V_\Pi:=\frac1m\sum_{k=1}^m
 H(X_{\Pi(2k-1)},X_{\Pi(2k)}).
 \label{eq:permutationaverage}
\end{equation}

\textbf{Step 1: recover the U-statistic by averaging pairings.}
A fixed unordered pair $\{i,j\}$ is selected in one of the $m$ slots with
probability $2m/[N(N-1)]$.  Its coefficient in
$\E_\Pi[V_\Pi\mid X_1,\ldots,X_N]$ is therefore
\begin{equation}
 \frac1m\frac{2m}{N(N-1)}=\binom N2^{-1}.
\end{equation}
Using \cref{eq:ustatunordered},
\begin{equation}
 \E_\Pi[V_\Pi\mid X_1,\ldots,X_N]=\widehat{\DT}_N.
 \label{eq:pairingidentity}
\end{equation}

\textbf{Step 2: bound the moment-generating function.}
For every $\lambda>0$, convexity of the exponential and
\cref{eq:pairingidentity} give
\begin{align}
 \E\exp\{\lambda(\widehat{\DT}_N-\theta)\}
 &\le\E\E_\Pi\exp\{\lambda(V_\Pi-\theta)\}.
 \label{eq:mgfjensen}
\end{align}
Conditional on a fixed permutation, the $m$ kernel values in
\cref{eq:permutationaverage} use disjoint independent sample pairs.  Each has
mean $\theta$ and lies in $[-M,M]$ because
$|A_i|\le M_U$ and $|B_i|\le M_P$.  Hoeffding's lemma (derived explicitly for a bounded interval in Step 3 of the proof of \cref{prop:onesidedwos}) therefore gives
\begin{equation}
 \E\exp\{\lambda(V_\Pi-\theta)\}
 \le\exp\left(\frac{\lambda^2M^2}{2m}\right).
 \label{eq:mgfhoeffding}
\end{equation}

\textbf{Step 3: Chernoff optimization.}
Markov's inequality and \cref{eq:mgfhoeffding} imply
\begin{equation}
 \Prob(\widehat{\DT}_N-\theta\ge t)
 \le\exp\left(-\lambda t+\frac{\lambda^2M^2}{2m}\right).
\end{equation}
The minimizing value is $\lambda=mt/M^2$.  Substitution produces
\begin{equation}
 \exp\!\left(-\frac{mt^2}{2M^2}\right).
\end{equation}
Applying the same argument to
$-\widehat{\DT}_N$ and adding the two tail probabilities proves
\cref{eq:hoeffdingtail}.

\textbf{Step 4: invert the tail and cover multiple candidates.}
Setting the right side of \cref{eq:hoeffdingtail} equal to $\delta$ gives
\cref{eq:hoeffdingradius}.  For candidate events $F_1,\ldots,F_K$, each with
failure probability at most $\delta/K$, Boole's inequality gives
$\Prob(\cup_kF_k)\le\sum_k\Prob(F_k)\le\delta$.  This argument never assumes
independence across candidate locations.
\end{proof}

\subsection{Marginal bias propagation}

\begin{proof}[Proof of \cref{prop:biaspropagation}]
Put $\bar u=\E U=u+e_U$ and $\bar p=\E P=p+e_P$, where
$|e_U|\le\beta_U$ and $|e_P|\le\beta_P$.  The same cross-index factorization
as in \cref{eq:crossmean} gives
\begin{equation}
 \E\widehat{\DT}_N=-(\bar u-h_z)\bar p.
\end{equation}
Subtracting $\DT=-(u-h_z)p$ and expanding yields the exact identity
\begin{equation}
 \E\widehat{\DT}_N-\DT
 =-p e_U-(u-h_z)e_P-e_Ue_P.
\end{equation}
The triangle inequality and the marginal bounds prove
\cref{eq:productbias}.
\end{proof}

\section{Finite-Difference Lower Bound and Complexity}
\label{app:rareevent}

\begin{proof}[Proof of \cref{thm:rarevariance}]
\textbf{Step 1: lower-bound the per-path second moment.}
By \cref{eq:rarezero},
\begin{align}
 \E Y_\eps^2
 &=\E[Y_\eps^2\mathbf 1_{A_\eps}]
 =\Prob(A_\eps)\E[Y_\eps^2\mid A_\eps]\\
 &\ge c_p c_2\rho_z(\eps).
 \label{eq:raresecondmoment}
\end{align}

\textbf{Step 2: control the squared mean.}
By \cref{eq:raremoments},
\begin{equation}
 (\E Y_\eps)^2\le C_1^2\rho_z(\eps)^2.
\end{equation}
If $\rho_z(\eps)\le c_pc_2/(2C_1^2)$, this is at most
$\tfrac12c_pc_2\rho_z(\eps)$.  Subtracting it from
\cref{eq:raresecondmoment} gives
\begin{equation}
 \Var(Y_\eps)\ge\frac12c_pc_2\rho_z(\eps).
 \label{eq:rareperpathvariance}
\end{equation}

\textbf{Step 3: scale the independent sample mean.}
Independence of the $N$ replicates gives
\begin{align}
 \Var(\widehat{\DT}^{\mathrm{FD}}_{N,\eps})
 &=\frac{1}{N\rho_z(\eps)^2}\Var(Y_\eps)\\
 &\ge\frac{c_pc_2}{2N\rho_z(\eps)},
\end{align}
which is \cref{eq:rarevariance}.
\end{proof}

\begin{proof}[Proof of \cref{cor:ranking}]
Apply \cref{eq:hoeffdingradius} to each candidate with failure probability
$\delta/K$.  The common radius is at most $\Delta$ precisely when
\begin{equation}
 M\sqrt{\frac{2\log(2K/\delta)}{\lfloor N/2\rfloor}}
 \le\Delta,
\end{equation}
which rearranges to \cref{eq:directcomplexity}.  The union bound in
\cref{thm:concentration} then gives simultaneous error at most $\Delta$.
For finite differencing, \cref{eq:rarevariance} and the requirement
$\Var(\widehat{\DT}^{\mathrm{FD}}_{N,\eps})\le\Delta^2$ imply
\cref{eq:fdcomplexity} by direct rearrangement.
\end{proof}

\section{Certification and Adaptive-Selection Proofs}
\label{app:certificate}

\begin{proof}[Proof of \cref{thm:certificate}]
Let
\begin{equation}
 E:=\{|\widehat{\DT}_N-\E\widehat{\DT}_N|
          \le r_N(\delta)\}.
\end{equation}
By hypothesis, $\Prob(E)\ge1-\delta$.

\textbf{Step 1: upper-bound the true derivative on $E$.}
The confidence event and \cref{eq:biasbudget} give
\begin{align}
 \DT(z)
 &\le\E\widehat{\DT}_N(z)
       +b_{\mathrm{term}}+b_{\mathrm{cache}}\\
 &\le\widehat{\DT}_N(z)+r_N(\delta)
       +b_{\mathrm{term}}+b_{\mathrm{cache}}.
 \label{eq:derivativeupperbound}
\end{align}

\textbf{Step 2: insert the deterministic remainder.}
Since $\rho_z(\eps)>0$, the upper side of
\cref{eq:remainderbound} and \cref{eq:derivativeupperbound} imply
\begin{align}
 \J(\Omega_{\eps,z})-\J(\Omega)
 &\le\rho_z(\eps)
       \{\DT(z)+C_z\rho_z(\eps)^\alpha\}\\
 &\le\rho_z(\eps)\{
 \widehat{\DT}_N+r_N+b_{\mathrm{term}}\notag\\[-2pt]
 &\hspace{23mm}{}+b_{\mathrm{cache}}
 +C_z\rho_z(\eps)^\alpha\}.
 \label{eq:certificatebound}
\end{align}
The braced term is negative by \cref{eq:certificate}; hence the objective
change is negative on $E$.  The event has probability at least $1-\delta$.
\end{proof}

\begin{proof}[Proof of \cref{cor:adaptiveradius}]
The event $E$ in the preceding proof is independent of which radius is chosen.
On $E$, \cref{eq:certificatebound} holds simultaneously for every
$\eps\in(0,\eps_0]$ because the remainder assumption is uniform over that
interval.  Therefore it also holds at the measurable random value
$\widehat\eps$.  If \cref{eq:certificate} holds there, the right side is
negative.  Thus failure can occur only on $E^c$, whose probability is at most
$\delta$.  Since one event covers the continuum of radii, no discretization or
union bound is required.
\end{proof}

\begin{proof}[Proof of \cref{prop:freshvalidation}]
Let $B_t$ be the event that round $t$ accepts a non-improving insertion.  The
candidate and radius are $\mathcal F_{t-1}$-measurable.  Conditional on that
history, the proof of \cref{thm:certificate} shows that an accepted insertion
is improving whenever $E_t$ occurs.  Hence
\begin{equation}
 B_t\subseteq E_t^c
 \quad\text{and}\quad
 \Prob(B_t\mid\mathcal F_{t-1})
 \le\Prob(E_t^c\mid\mathcal F_{t-1})\le\delta_t.
\end{equation}
Taking expectations and using the tower property gives
$\Prob(B_t)\le\delta_t$.  Finally, countable subadditivity yields
\begin{equation}
 \Prob\left(\bigcup_{t\ge1}B_t\right)
 \le\sum_{t\ge1}\Prob(B_t)
 \le\sum_{t\ge1}\delta_t,
\end{equation}
which is \cref{eq:sequentialfamilywise}.  No independence between validation
batches is used; only the conditional guarantee is required.
\end{proof}

\section{Source Integration and Connected Refinement}
\label{app:channelimplementation}

\subsection{Poisson source integration}

For completeness, the inhomogeneous three-dimensional Poisson implementation
used in our engineering benchmark samples the Green source term without a
volume grid.  The following identity also fixes the normalization used in our
implementation.

\begin{proposition}[Unbiased source sample on a three-dimensional WoS ball]
\label{prop:greenballsample}
Let $B_R(x)\subset\Omega$, let $\Theta$ be uniform on $\mathbb S^2$, and let
$T$ be independent of $\Theta$ with density $p_T(t)=6t(1-t)$ on $(0,1)$.
For every measurable source $f$ with $\int_{B_R(x)}G_{B_R(x)}(x,y)|f(y)|\,dy<\infty$,
\begin{equation}
 \int_{B_R(x)}G_{B_R(x)}(x,y)f(y)\,dy
 =\frac{R^2}{6}\,\E[f(x+RT\Theta)],
 \label{eq:greenballsample}
\end{equation}
where $G_{B_R(x)}(x,y)=(4\pi)^{-1}
(\|y-x\|^{-1}-R^{-1})$ is the Dirichlet Green kernel evaluated at the ball
center.
\end{proposition}

\begin{proof}
Denote the left-hand side of \cref{eq:greenballsample} by $I_R(x)$ and write
$y=x+r\theta$.  Spherical coordinates, followed by $r=Rt$, give
\begin{align*}
 I_R(x)
 &=\frac1{4\pi}\int_{\mathbb S^2}\int_0^R
   \left(r-\frac{r^2}{R}\right)f(x+r\theta)\,dr\,d\theta\\
 &=\frac{R^2}{4\pi}\int_{\mathbb S^2}\int_0^1
   t(1-t)f(x+Rt\theta)\,dt\,d\theta.
\end{align*}
Uniform $\Theta$ has surface density $(4\pi)^{-1}$ and $T$ has density
$6t(1-t)$; substituting both densities into the expectation on the right of
\cref{eq:greenballsample} gives exactly the last integral.  Hence one draw of
$(T,\Theta)$ yields an unbiased contribution $R^2f(x+RT\Theta)/6$.
\end{proof}

\subsection{Connected-channel parameterization and refinement}

Topology creation need not end with disconnected cavities.  Let
$q=(q_0,\ldots,q_m)\in(\R^3)^{m+1}$ be control points and let a fixed
subdivision matrix $W\in\R^{L\times(m+1)}$ have nonnegative entries and row
sums one.  Dense centers, the removed channel, and the retained domain are
\begin{equation}
 \begin{aligned}
  s_i(q)&=\sum_{j=0}^{m}W_{ij}q_j,\\
  H(q,r)&=\bigcup_{i=1}^{L}\overline B_r(s_i(q)),
  &\Omega(q,r)&=\Omega\setminus H(q,r).
 \end{aligned}
 \label{eq:channelparameterization}
\end{equation}
Every accepted trial must satisfy $\|s_{i+1}-s_i\|<2r$ for all consecutive
centers, checked again after radius recalibration.  Thus the interiors of
successive balls overlap, making both the removed union and its open interior
path connected.  The
topology stage selects two separated terminals and connects them through a
topological-score-weighted candidate graph; Differential WoS then moves the
interior controls while the endpoints remain fixed.
The current housing experiment keeps every ball strictly inside the housing;
its terminals are interior design points.  It therefore demonstrates an
internally connected cavity, not an inlet-to-outlet conduit or a guarantee of
powder evacuation.  In particular, it would remain a void island under the
exterior-reachability definition of Andrade et al.~\cite{andrade2026connectivity}.
The two constraints must not be treated as interchangeable baselines.

Let $\Sigma$ be the union of sphere-intersection seams on $\partial H$.  For
an exposed patch $\Gamma_i$ and offset $\ell>0$, define the admissible part
of that patch by
\begin{equation}
 \Gamma_{i,\ell}^{\mathrm{safe}}=
 \left\{x\in\Gamma_i:
 \operatorname{dist}(x,\Sigma)>\ell,\;
 (x,x-\ell n(x)]\subset\Omega(q,r)\right\}.
 \label{eq:channelsafepatch}
\end{equation}
The Euclidean seam distance is used here; replacing the unit strip width by
any fixed positive multiple changes only constants.

\begin{proposition}[Differential WoS center derivative]
\label{prop:channel-diffwos}
Fix $r>0$ and a regular channel configuration for which the exposed spherical
patches $\Gamma_i\subset\partial H(q,r)$ are $C^2$ away from $\Sigma$.
Assume the piecewise-smooth solution map is shape differentiable under
independent center translations.  Moreover, for every compact probe set
$K\Subset\Omega(q,r)$, assume there are $C_K,\ell_K>0$ such that, for
$0<\ell\le\ell_K$ and $\xi=X_\tau$,
\begin{align}
 \sup_{y\in K}\E_y\!\left[
  \mathbf 1_{\{I(\xi)=i,\,\xi\notin\Gamma_{i,\ell}^{\mathrm{safe}}\}}
  |\partial_nu(\xi)|\right]&\le C_K\ell,
 \label{eq:weightedseamstrip}\\
 \sup_{x\in\Gamma_{i,\ell}^{\mathrm{safe}}}
 \left|\frac{u(x-\ell n(x))}{\ell}+\partial_nu(x)\right|
 &\le C_K\ell.
 \label{eq:safenormaltaylor}
\end{align}
Let $u$ solve
$-\Delta u=f$ in $\Omega(q,r)$ with
zero Dirichlet data.  For $y\in\Omega(q,r)$, let $X_\tau$ be Brownian exit,
$I(X_\tau)$ its exposed-ball index, and $n$ the outward normal of the retained
domain.  If $U_x$ is an unbiased independent WoS estimate of $u(x)$, then
\begin{equation}
 G_{i,\ell}(y)=
 \mathbf 1_{\{I(X_\tau)=i,\,
 X_\tau\in\Gamma_{i,\ell}^{\mathrm{safe}}\}}
 n(X_\tau)\frac{U_{X_\tau-\ell n(X_\tau)}}{\ell}
 \label{eq:diffwoscenterestimator}
\end{equation}
satisfies, uniformly on compact probe sets,
\begin{equation}
 \E G_{i,\ell}(y)=\nabla_{s_i}u(y)+O(\ell)
 \quad\text{as }\ell\downarrow0.
 \label{eq:diffwoscenterbias}
\end{equation}
For any fixed finite set of interior probes, the weighted estimate has $O(\ell)$ bias for the corresponding weighted derivative sum. To assert this rate for compliance $J(q,r)=\int_{\Omega(q,r)}fu\,dx$, additionally assume that differentiation under the domain integral is valid and that, for all sufficiently small $\ell$, the pointwise estimator error is bounded by $\ell C(y)$ with $\int_{\Omega(q,r)}|f(y)|C(y)\,dy<\infty$. A finite probe rule then adds its separate quadrature error.
\end{proposition}

\begin{proposition}[Control chain rule and length tangency]
\label{prop:channel-chainrule}
Suppose $J$ is differentiable with respect to the dense centers at a regular
configuration, and write $G_i=\nabla_{s_i}J$.  Then
\begin{equation}
 \nabla_{q_j}J=\sum_{i=1}^{L}W_{ij}G_i.
 \label{eq:channelchainrule}
\end{equation}
For the centerline length
$\mathcal L(q)=\sum_{j=0}^{m-1}\|q_{j+1}-q_j\|$, let $c$ be its gradient
with the endpoint blocks set to zero.  Given any endpoint-fixed proposal $a$,
define
\begin{equation}
 d=a-c\frac{\langle c,a\rangle}{\|c\|^2}
 \quad(c\ne0),
 \label{eq:lengthprojection}
\end{equation}
and $d=a$ if $c=0$.  Then the endpoints do not move and
$D\mathcal L(q)[d]=0$.  In particular, let $g$ be the objective gradient with
its endpoint blocks set to zero.  Choosing $a=-g$ gives
$DJ(q)[d]=-\|P_{c^\perp}g\|^2\le0$.  Our implementation first
smooths $a$ and therefore relies on fresh line-search validation, rather than
this last inequality, for accepted descent.
\end{proposition}

The radius of every trial is recalibrated on one fixed Sobol set.  If
$Y_1,\ldots,Y_M$ fill a box of volume $V_B$, define
\begin{equation}
 \widehat V_M(q,r)=\frac{V_B}{M}\sum_{k=1}^{M}
 \mathbf 1_{\{Y_k\in\Omega,\;\min_i\|Y_k-s_i(q)\|\le r\}}.
 \label{eq:qmcchannelvolume}
\end{equation}

\begin{proposition}[Discrete equal-volume calibration]
\label{prop:qmc-volume}
Let $k$ be the nearest integer to $MV_\star/V_B$.  If the $k$th and
$(k+1)$st eligible center distances are distinct, choosing $r$ strictly
between them gives exactly $k$ removed samples and
\begin{equation}
 |\widehat V_M(q,r)-V_\star|\le \frac{V_B}{2M}.
 \label{eq:qmcvolumequantum}
\end{equation}
Thus all compared connected channels have the same discrete volume, up to one
predeclared rounding quantum; this is a statement about the fixed QMC measure,
not an exact analytic union-volume formula.
\end{proposition}

\begin{proposition}[Fresh paired shape acceptance]
\label{prop:fresh-shape}
At shape round $t$, condition on the complete history $\mathcal F_{t-1}$,
including the Differential WoS direction and finitely many trial scales.  Let
$\Delta_{t,k}=J(q_{t,k},r_{t,k})-J(q_t,r_t)$, where every $r_{t,k}$ obeys
\cref{eq:qmcchannelvolume}.  Suppose simultaneous upper bounds satisfy
\begin{equation}
 \Prob\!\left(\Delta_{t,k}\le U_{t,k}\ \text{for every }k
 \mid\mathcal F_{t-1}\right)\ge1-\delta_t.
 \label{eq:shapesimultaneouscoverage}
\end{equation}
If a trial is accepted only when its $U_{t,k}<0$, then
\begin{equation}
 \Prob(\text{some accepted shape step is non-improving})
 \le\sum_t\delta_t.
 \label{eq:shapefamilywise}
\end{equation}
The conclusion is independent of bias in the proposed direction: coverage of
the fresh objective difference, not gradient accuracy, certifies acceptance.
\end{proposition}



\section{Connected-Channel Refinement Proofs}
\label{app:channelproofs}

\begin{proof}[Proof of \cref{prop:channel-diffwos}]
We prove the vector identity by testing it against an arbitrary translation
$v\in\R^3$ of ball $i$.

\textbf{Step 1: derive the boundary condition for the shape derivative.}
On the interior of the exposed patch $\Gamma_i$, translation of center $i$
has velocity $V=v$; every other open spherical patch and the outer boundary
have $V=0$.  The proposition's shape-differentiability assumption permits
patchwise trace differentiation; the seam itself is excluded.  Write $u_t$
for the perturbed solution and $\dot u$ for its Eulerian shape derivative at
$t=0$.  Differentiating the
identity $u_t(T_t(x))=0$ for $x\in\Gamma_i$ gives
\begin{equation}
 \dot u(x)+\nabla u(x)\cdot V(x)=0.
 \label{eq:channelmaterialboundary}
\end{equation}
The trace of $u$ is constant on $\Gamma_i$, so its tangential gradient
vanishes there.  Hence $\nabla u=(\partial_nu)n$ and
\begin{equation}
 \dot u=-(n\cdot v)\partial_nu\quad\text{on }\Gamma_i.
 \label{eq:channelshapecondition}
\end{equation}
The source $f$ and the outer boundary are fixed.  Differentiating the weak
Poisson equation with test functions compactly supported in the current
domain therefore gives
\begin{equation}
 -\Delta\dot u=0\quad\text{in }\Omega(q,r),
 \label{eq:channelharmonicderivative}
\end{equation}
with zero trace on all stationary Dirichlet patches.  Equations
\cref{eq:channelshapecondition,eq:channelharmonicderivative} are the standard
differential boundary-value problem, here obtained directly from the moving
trace.

\textbf{Step 2: approximate the normal derivative on the safe patch.}
For $x\in\Gamma_{i,\ell}^{\mathrm{safe}}$, the segment in
\cref{eq:channelsafepatch} lies in the retained domain.  Whenever a uniform
$C^2$ extension is available, Taylor's theorem along this segment and
$u(x)=0$ give
\begin{align}
 u(x-\ell n)
 &=u(x)-\ell\partial_nu(x)
   +\frac{\ell^2}{2}
     n^T\nabla^2u(x-\theta\ell n)n \\
 &=-\ell\partial_nu(x)+O(\ell^2),
 \label{eq:backwardnormaltaylor}
\end{align}
for some $\theta\in(0,1)$.  Uniform $C^2$ bounds on the regular exposed
patches imply
\begin{equation}
 -\partial_nu(x)=\frac{u(x-\ell n(x))}{\ell}+O(\ell)
 \label{eq:backwardnormalderivative}
\end{equation}
there.  In the proposition we use the explicitly stated uniform version
\cref{eq:safenormaltaylor}; this separates the needed analytic estimate from
the merely qualitative fact that $\Sigma$ has zero harmonic measure.
Combining
\cref{eq:channelshapecondition,eq:backwardnormalderivative} yields
\begin{equation}
 \dot u(x)=(n(x)\cdot v)
 \frac{u(x-\ell n(x))}{\ell}+O(\ell\|v\|).
 \label{eq:offsetboundarydata}
\end{equation}

\textbf{Step 3: truncate the harmonic exit representation.}
Let $X_\tau$ denote Brownian exit from $\Omega(q,r)$, and abbreviate
$\xi=X_\tau$ and $n_\tau=n(X_\tau)$.  The harmonic representation of
\cref{eq:channelharmonicderivative} first gives the exact identity
\begin{align}
 D_{s_i}u(y)[v]
 &=\E_y\!\left[
   \mathbf 1_{\{I(\xi)=i\}}
   \bigl(-(n_\tau\cdot v)\partial_nu(\xi)\bigr)\right].
 \label{eq:exactharmonicshaperepresentation}
\end{align}
Write $S_{i,\ell}=\{I(\xi)=i,\,
\xi\in\Gamma_{i,\ell}^{\mathrm{safe}}\}$.  On $S_{i,\ell}$,
\cref{eq:safenormaltaylor} bounds the error made by replacing
$-\partial_nu(\xi)$ with $u(\xi-\ell n_\tau)/\ell$ by $C_K\ell$.
Because harmonic measure has total mass at most one, its contribution is at
most $C_K\ell\|v\|$.  On $\{I(\xi)=i\}\setminus S_{i,\ell}$,
\cref{eq:weightedseamstrip} bounds the discarded part of
\cref{eq:exactharmonicshaperepresentation} by the same order.  Therefore
\begin{align}
 D_{s_i}u(y)[v]
 &=\E_y\!\left[
   \mathbf 1_{S_{i,\ell}}(n_\tau\cdot v)
   \frac{u(\xi-\ell n_\tau)}{\ell}\right]
   +O_K(\ell\|v\|).
 \label{eq:harmonicshaperepresentation}
\end{align}
This is the precise point at which the quantitative strip assumption, rather
than only measure zero of the limiting seam, is required.

\textbf{Step 4: insert the nested WoS estimate.}
Condition on $X_\tau$ on the safe event $S_{i,\ell}$.  By hypothesis,
\begin{equation}
 \E[U_{X_\tau-\ell n(X_\tau)}\mid X_\tau]
 =u(X_\tau-\ell n(X_\tau)).
\end{equation}
Taking iterated expectations in \cref{eq:diffwoscenterestimator} and comparing
with \cref{eq:harmonicshaperepresentation} shows that, for every $v$,
\begin{equation}
 v\cdot\E G_{i,\ell}(y)
 =D_{s_i}u(y)[v]+O(\ell\|v\|).
\end{equation}
Duality in the finite-dimensional space $\R^3$ proves
\cref{eq:diffwoscenterbias}.

\textbf{Step 5: differentiate compliance.}
Reynolds' transport formula gives
\begin{equation}
 D_{s_i}J[v]
 =\int_{\Omega(q,r)}fD_{s_i}u[v]\,dx
  +\int_{\partial\Omega(q,r)}fu(V\cdot n)\,dS.
 \label{eq:channelreynolds}
\end{equation}
The boundary integral vanishes because $u=0$ on every moving and fixed
Dirichlet patch. For a fixed finite interior probe set, compact-set uniformity in \cref{eq:diffwoscenterbias} directly bounds the weighted sum of estimator errors by $O(\ell)$. For the continuum objective, the additional integrable-envelope hypothesis gives
\[
 \begin{aligned}
 &\left\|\int_{\Omega(q,r)}f(y)
 \bigl(\E G_{i,\ell}(y)-\nabla_{s_i}u(y)\bigr)\,dy\right\|\\
 &\qquad\le \ell\int_{\Omega(q,r)}|f(y)|C(y)\,dy.
 \end{aligned}
\]
This proves the continuum rate under the added hypothesis. Compact-set bounds alone would not control probes approaching the boundary. Replacing the integral by a fixed weighted rule adds the separate quadrature error.
\end{proof}

\begin{proof}[Proof of \cref{prop:channel-chainrule}]
\textbf{Step 1: differentiate the dense centers.}
For a control perturbation $\dot q=(\dot q_0,\ldots,\dot q_m)$,
linearity of \cref{eq:channelparameterization} gives
\begin{equation}
 \dot s_i=\sum_{j=0}^{m}W_{ij}\dot q_j.
 \label{eq:densecenterdifferential}
\end{equation}

\textbf{Step 2: apply the ordinary finite-dimensional chain rule.}
At the assumed differentiability point,
\begin{align}
 DJ(q)[\dot q]
 &=\sum_{i=1}^{L}G_i\cdot\dot s_i
 =\sum_{i=1}^{L}G_i\cdot
   \sum_{j=0}^{m}W_{ij}\dot q_j\\
 &=\sum_{j=0}^{m}
   \left(\sum_{i=1}^{L}W_{ij}G_i\right)\cdot\dot q_j.
 \label{eq:channelchainproof}
\end{align}
Since this holds for every $\dot q$, comparison with the definition of the
Euclidean gradient proves \cref{eq:channelchainrule}.

\textbf{Step 3: compute the length gradient.}
Put
$e_j=(q_{j+1}-q_j)/\|q_{j+1}-q_j\|$ and assume no consecutive controls
coincide.  Direct differentiation yields
\begin{equation}
 \begin{aligned}
  \nabla_{q_0}\mathcal L&=-e_0,\\
  \nabla_{q_j}\mathcal L&=e_{j-1}-e_j &&(1\le j<m),\\
  \nabla_{q_m}\mathcal L&=e_{m-1}.
 \end{aligned}
 \label{eq:lengthgradientexplicit}
\end{equation}
For endpoint-fixed directions only the interior blocks contribute, so replacing
the two endpoint blocks by zero gives the vector $c$ in the proposition and
$D\mathcal L(q)[d]=\langle c,d\rangle$.

\textbf{Step 4: verify tangency.}
When $c\ne0$, substitution of \cref{eq:lengthprojection} gives
\begin{equation}
 \langle c,d\rangle
 =\langle c,a\rangle
  -\|c\|^2\frac{\langle c,a\rangle}{\|c\|^2}=0.
\end{equation}
When $c=0$, the same conclusion is immediate.  Both $a$ and $c$ have zero
endpoint blocks, hence so does $d$.

\textbf{Step 5: verify descent for the unsmoothed projected gradient.}
Let $g$ denote $\nabla_qJ$ with its endpoint blocks set to zero and let
$P_{c^\perp}$ be orthogonal projection within that endpoint-fixed subspace.
Choosing $a=-g$ makes $d=-P_{c^\perp}g$.  Therefore
\begin{equation}
 DJ(q)[d]=\langle g,d\rangle
 =-\langle g,P_{c^\perp}g\rangle
 =-\|P_{c^\perp}g\|^2\le0.
\end{equation}
The final equality uses symmetry and idempotence of an orthogonal projection.
No descent assertion follows after an arbitrary smoothing operator; this is
why the implemented smoothed proposal is subjected to fresh evaluation.
\end{proof}

\begin{proof}[Proof of \cref{prop:qmc-volume}]
For each sample define
\begin{equation}
 d_j=\begin{cases}
 \min_i\|Y_j-s_i(q)\|,&Y_j\in\Omega,\\
 +\infty,&Y_j\notin\Omega.
 \end{cases}
\end{equation}
Let $d_{(1)}\le\cdots\le d_{(M)}$ be the order statistics.  If
$d_{(k)}<r<d_{(k+1)}$, exactly the first $k$ indicators in
\cref{eq:qmcchannelvolume} equal one.  Hence
\begin{equation}
 \widehat V_M(q,r)=\frac{kV_B}{M}.
\end{equation}
Since $k$ is a nearest integer to $MV_\star/V_B$,
\begin{equation}
 \left|k-\frac{MV_\star}{V_B}\right|\le\frac12.
\end{equation}
Multiplication by $V_B/M$ proves \cref{eq:qmcvolumequantum}.  The distinctness
assumption prevents ambiguity at the threshold.  If ties occur, a deterministic
tie rule retains a one-quantum bound but need not attain exactly $k$ samples.
\end{proof}

\begin{proof}[Proof of \cref{prop:fresh-shape}]
Let $E_t$ be the simultaneous-coverage event inside
\cref{eq:shapesimultaneouscoverage}, and let $B_t$ be the event that round $t$
accepts a non-improving trial.  On $E_t$, every accepted index $k$ satisfies
\begin{equation}
 \Delta_{t,k}\le U_{t,k}<0.
\end{equation}
Consequently $B_t\subseteq E_t^c$.  Conditional coverage and the tower
property imply
\begin{align}
 \Prob(B_t)
 &=\E[\Prob(B_t\mid\mathcal F_{t-1})]\\
 &\le\E[\Prob(E_t^c\mid\mathcal F_{t-1})]
 \le\delta_t.
\end{align}
Finally, countable subadditivity gives
\begin{equation}
 \Prob\!\left(\bigcup_tB_t\right)
 \le\sum_t\Prob(B_t)\le\sum_t\delta_t,
\end{equation}
which is \cref{eq:shapefamilywise}.  The direction, tested scales, common
random numbers within each paired difference, and deterministic radius
recalibration are all part of the conditioning history.  No correctness
assumption about the Differential WoS proposal is used after that conditioning.
\end{proof}

\section{Reproducibility and Interpretation}
\label{app:checklist}
\label{app:claims}

Detailed protocols appear in \cref{app:experiments,sec:critical}; the count--radius study retains all 72 designs and separates design-seed ranges from conditional evaluation uncertainty. Rendering cuts and scalar smoothing affect only visualization, not the physical domain or reported objectives.

The deterministic capacity expansion uses the analytic hypotheses stated in the main text. Cross-fitted unbiasedness assumes exact marginal oracle means; implemented termination error remains separate. The finite-insertion and shape-acceptance results control erroneous acceptances under their coverage premises, rather than giving posterior probabilities conditional on an accepted design. No application result is presented as a proof that its finite-size, oracle-bias, or interval assumptions hold.
